\section{Related Work}
\label{sec:related}

\begin{table*}[t]
\caption{Positioning. ``Evidence signed by'' asks who signs the record of
an execution. ``Resume'' asks whether a task can continue on a different
provider from a state verified by the requester. ``n/s'': not specified
in the cited source; ``---'': out of the system's scope, or, in the
Evaluation column, no evaluation reported.}
\label{tab:related}
\small
\begin{tabular}{@{}lllllp{3.2cm}@{}}
\toprule
System & Evidence signed by & Trust computed & Resume & Metering $\leftrightarrow$ evidence & Evaluation \\
\midrule
Agentic P2P~\cite{wang2026agenticp2p} & provider & local + federated ledger$^{a}$ & no & no & simulation \\
MG-EigenTrust~\cite{zhang2026distributed} & n/s & global (EigenTrust-style) & no & no & simulation, prototype \\
ERC-8004~\cite{erc8004} & feedback author & global, on-chain & no & n/s & measured by~\cite{xiong2026erc8004} \\
INAM~\cite{inam} & worker, then requester & registry (a hint)$^{b}$ & no & no$^{c}$ & live deployment$^{d}$ \\
Cairn~\cite{chard2026cairn} & n/s & community (decayed Beta) & no & no & simulation, case study \\
AgentReputation~\cite{chishti2026agentreputation} & n/s & decentralized service & no & no & vision paper \\
\midrule
Signet~\cite{signet} & agent (request), server (response) & --- & no & no & --- \\
Sello~\cite{figuera2026sello} & receiver & --- & no & no & microbenchmarks \\
ACTA~\cite{acta2026} & decision maker & --- & no & no & --- \\
Agent Passport~\cite{aps} & agent & --- & no & no & --- \\
SCITT (RFC 9943)~\cite{rfc9943} & single issuer & --- & --- & --- & --- \\
PeerReview~\cite{haeberlen2007peerreview} & each node's log & --- & --- & --- & implementation \\
TrustChain~\cite{otte2020trustchain} & both parties & per agent (NetFlow) & --- & --- & implementation \\
\midrule
A2A~\cite{a2a2026} & (signed Agent Card) & --- & no & --- & --- \\
ADS~\cite{ads2026}, NANDA~\cite{nanda2025} & (signed records) & --- & --- & --- & --- \\
MPP~\cite{mpp2026}, x402 V2~\cite{x402v2} & payer (vouchers) & --- & --- & usage, not execution & --- \\
AP2~\cite{ap2} & user (mandates) & --- & --- & authorization & --- \\
\midrule
EigenTrust~\cite{kamvar2003eigentrust} & --- & global & --- & --- & simulation \\
Credence~\cite{walsh2006credence} & voter & per client (correlation, flow) & --- & --- & deployment \\
Beta reputation~\cite{josang2002beta} & --- & aggregated ratings & --- & --- & --- \\
\midrule
\textbf{\sys{}} & \textbf{requester and provider} & \textbf{observer-local} & \textbf{yes} & \textbf{voucher digest} & \textbf{pre-registered sim.; 2 impls} \\
\bottomrule
\end{tabular}

\smallskip
\raggedright\footnotesize $^{a}$ As described in the paper's prose. Its
released simulator ranks providers by a hint the provider reports about
itself (\S\ref{sec:background}).
$^{b}$ The registry computes one score per agent; the specification calls
it a hint that relying parties re-derive from the signed receipts.
$^{c}$ A settlement reference reported by the parties and not verified.
$^{d}$ A live registry. We found no published measurement of its scoring
under attack (checked 2026-11-13).
\end{table*}

\paragraph{Agent networks and reputation.}
Wang et al.~\cite{wang2026agenticp2p} framed agent delegation as the
successor of P2P content distribution. We adopt that framing and
replicate their evaluation. We differ in who signs the evidence, which
receipts are admissible, and how issuers are weighted; like the ledger
they propose, our trust is local to each observer. We also add
resumption. Their group's follow-up,
MG-EigenTrust, proposes multi-topic EigenTrust against cross-topic
collusion~\cite{zhang2026distributed}. Our E1 result for W2 is a warning
for any propagation-based scheme: once camouflaging Sybils outnumber
honest raters, propagation amplifies them.

ERC-8004~\cite{erc8004} provides on-chain identity, reputation and
validation registries. An empirical study found its reputation registry
unusable as a trust signal, with 59--91\% of reviewers showing coordinated
Sybil behaviour~\cite{xiong2026erc8004}. E1 does not reproduce that finding
directly. Its global ledger B1 holds up under farming (0.094 at
$\alpha{=}0.2$, half of random) and collapses under slander. Unlike
ERC-8004's feedback, it also admits positive evidence only when
dual-signed.

INAM~\cite{inam} is the closest agent-specific system, so we compare it
precisely. As of its specification version 0.44-draft:
\begin{itemize}
\item \textbf{Receipts.} The signing order is \sys{}'s over A2A: the
  worker drafts an execution receipt and the requester countersigns it.
  Only a receipt both have signed counts toward reputation, as under R-4.
  INAM identities are did:key Ed25519 keys, byte-compatible with
  AgentIDs. An INAM receipt covers one interaction and commits to a hash
  of the task specification and a hash of the output. A \sys{} task has
  countersigned checkpoints over a Merkle root of its streamed output.
  INAM's HTTP binding states that a streamed response does not fit it.
\item \textbf{Where trust is computed.} An INAM registry computes one
  score per agent for every caller. It weights each receipt by the
  counterparty's own score (a one-step relaxation of EigenTrust), caps
  repeated receipts from one counterparty, and gives no weight to a
  counterparty with no receipt outside the agent's own set of
  counterparties. The specification calls this score a hint and tells a
  relying party to re-derive the evidence from the signed receipts. It
  does not say how that party should weight what it re-derives;
  \sys{}'s projection is one answer. E1 did not test INAM's formula, but
  two of its findings concern similar mechanisms. First, the global
  ledger that admits unilateral complaints collapsed under slander; INAM
  adds a job poster's one-sided non-performance report to the worker's
  reputation, weighted by the poster's score. Second, propagation
  collapsed once camouflaging Sybils outnumbered honest issuers; INAM's
  specification notes that its one-hop check misses a ring whose members
  each obtain one outside receipt.
\item \textbf{What each lacks.} INAM has a transparency log of finalized
  receipts in the style of RFC~6962, verifiers authorized by the registry
  operator, dispute windows, key succession, and links to ERC-8004 and
  to x402, where a payer can refuse a payee whose INAM identity has not
  proved control of the receiving wallet. \sys{} has none of these. INAM
  has no resumption and no zones, and its settlement field is a reference
  the parties report, which INAM does not verify; it binds no usage to
  evidence.
\end{itemize}
The two compose. One key can serve as both identities, a \sys{} task can
end in an INAM receipt that names the task's TaskID, and a \sys{}
observer could admit INAM receipts as evidence under T-3. We have not
implemented either.

Cairn~\cite{chard2026cairn} aggregates agents' ratings of shared
resources with a time-decayed Beta model and tests it against lying,
collusion and camouflage. Its aggregation is community-wide. \sys{}'s is
per observer, and \sys{} admits positive evidence only if it is
countersigned. AgentReputation~\cite{chishti2026agentreputation} sketches
context-conditioned reputation with explicit verification regimes, which
is complementary to our admissibility rules.

\paragraph{Receipts and transparency.}
Signed receipts for agent actions are an active area. ACTA signs
access-control decisions~\cite{acta2026}. Agent Passport signs a receipt
for every delegated action~\cite{aps}. Sello inverts the trust boundary:
the receiving service signs what it observed, encrypts it to the agent's
owner, and publishes it to a witness-cosigned log~\cite{figuera2026sello}.
None of these defines how receipts feed trust, how a task resumes
elsewhere, or how usage binds to payment.

Signet~\cite{signet} is the closest agent-specific analogue of \sys{}'s
dual signatures, so we compare them precisely. As of Signet version 0.11:
\begin{itemize}
\item \textbf{Who attests the outcome.} In a Signet bilateral receipt the
  calling agent signs the request (tool, parameters, time). The server
  then co-signs its response and the outcome state (executed, failed,
  rejected, or requires approval). The outcome is therefore attested by
  the party that did the work. In \sys{} the provider proposes the
  outcome and usage, and the requester, the party that received the
  work, countersigns them. Only that countersignature makes a receipt
  positive evidence (R-4).
\item \textbf{Granularity.} Signet signs one tool call, request and
  response. \sys{} signs a task, with countersigned checkpoints over a
  Merkle root of its streamed output, which is what lets the task resume
  on another provider.
\item \textbf{What the evidence is for.} Signet targets audit and
  authorization: hash-chained local logs, delegation chains whose scope
  only narrows, policy attestations, revocation and replay protection.
  \sys{} has none of these. It retired delegation, and leaves key
  rotation and revocation out of scope. \sys{} instead defines how
  receipts become trust, through admissibility and observer-local
  weighting, and it binds usage to payment vouchers. Signet's spend
  limits constrain what an agent is authorized to do; they are not
  metering evidence.
\end{itemize}
The two are complementary. A Signet receipt could serve as the evidence
for a single tool call made inside a \sys{} task.

Outside the agent setting, the closest prior design is
TrustChain~\cite{otte2020trustchain}.
It records each interaction only with both parties' signatures, so a
Sybil can create trust only with itself, and each agent ranks others with
a flow computation (NetFlow) over the interaction graph. \sys{}'s R-4 and
T-3 make the same argument. \sys{} differs in three ways:
\begin{itemize}
\item its receipts are standalone artifacts, not entries in per-agent
  hash chains, so it has no chain-consistency or double-spend detection;
\item it weights issuers by agreement with the observer's own experience,
  not by contribution flow;
\item it adds countersigned mid-task checkpoints, resumption on another
  provider, and a binding to payment vouchers.
\end{itemize}

\paragraph{Attesting execution.}
A receipt records what two parties agreed happened. It does not prove
that a computation ran correctly. Stronger evidence is possible, at a
cost:
\begin{itemize}
\item \textbf{Verified execution.} Trusted-execution attestation, proofs
  of computation, and re-execution by validators are the options ERC-8004's
  validation registry anticipates~\cite{erc8004}.
\item \textbf{Accountability primitives.} A2M~\cite{chun2007a2m} and
  TrInc~\cite{levin2009trinc} use small trusted components to stop a
  party from telling different stories to different peers.
  AVM~\cite{haeberlen2010avm} records execution for audit by replay.
\end{itemize}
\sys{} needs none of these, and a receipt could be accompanied by any of
them. We chose receipts because they need neither trusted hardware nor
deterministic execution, and agent tasks rarely offer the latter.

SCITT~\cite{rfc9943} and Certificate Transparency~\cite{rfc6962} make
single-issuer statements auditable through transparency logs. \sys{}'s
receipts could be registered with such a log, but \sys{} does not need
one: evidence is held by the parties and disclosed at their discretion.
PeerReview~\cite{haeberlen2007peerreview} established accountability
through signed, tamper-evident logs with deterministic replay.
Countersigned checkpoints are a lightweight, application-level analogue,
for work that is not deterministic.

\paragraph{Payments.}
MPP sessions~\cite{mpp2026} and x402~V2~\cite{x402v2} already provide
payer-signed cumulative metering. These, like payment
channels~\cite{poon2016lightning}, prove that value was authorized, not
which work it bought. AP2's mandates~\cite{ap2} prove what the user
authorized. \sys{} adds only the reverse binding: a voucher digest inside
dual-signed execution evidence. Countersigning and paying per checkpoint
turn one exchange into many small ones, so either party's loss is at
most one step. Fair-exchange protocols obtain strong fairness instead,
with an optional trusted third party~\cite{asokan1997fairexchange}. \sys{}
accepts one step of exposure rather than introduce one.

\paragraph{Classic reputation and Sybil defence.}
Our scoring is Beta reputation~\cite{josang2002beta}, and W2 is a
personalized variant of EigenTrust~\cite{kamvar2003eigentrust}. Neither
weighting is new:
\begin{itemize}
\item \textbf{W1} is the credibility weighting of PeerTrust's personalized
  similarity measure~\cite{xiong2004peertrust} and of
  Credence~\cite{walsh2006credence}. Both weigh another rater by how
  closely its past ratings match one's own on commonly rated items.
\item \textbf{W2}'s analogue is Credence's second step: it combines those
  pairwise weights with a PageRank-style flow computation, from each
  client's own perspective.
\end{itemize}
Credence bounds Sybils with certificates from a central authority, an
external cost that \sys{} does not have. What E1 adds is a measurement in
the agent setting: propagation fails once camouflaging Sybils outnumber
honest raters. Direct agreement holds at 1.6 Sybils per honest requester and falls behind self-only trust by 6.4.
DSybil~\cite{yu2009dsybil} goes further. It gives provable, optimal
bounds on a user's loss against an unlimited number of Sybil identities,
under an assumption about how honest users vote. \sys{} has no comparable
guarantee (\S\ref{sec:security}). The
attack taxonomy follows~\cite{hoffman2009survey}, and whitewashing
follows Feldman et al.~\cite{feldman2006whitewashing}. Graph-based Sybil
defences such as SybilGuard~\cite{yu2006sybilguard} assume that
attack edges to honest nodes are scarce. Camouflage in E1 violates that
assumption by construction, and W2's collapse is the measured
consequence. S/Kademlia~\cite{baumgart2007skademlia} motivates D-9, and
E3c replicates its effect.
